\chapter{Parte 1: navegacion por campos potenciales}

La primera parte consiste en hacer que el robot avance hacia un objetivo. El objetivo puede ser un objeto \textit{mannequin}, Bill u otro dummy definido en la escena. La idea es calcular la posicion relativa entre robot y objetivo y convertirla en velocidad de avance y giro.

\section{Ley de control}

Sea \((x_r,y_r)\) la posicion del robot y \((x_g,y_g)\) la posicion del objetivo. El vector de atraccion se define como:

\[
\vec{F}_{att} = k_{att}\begin{bmatrix}x_g-x_r\\y_g-y_r\end{bmatrix}
\]

El angulo deseado se obtiene con \(\operatorname{atan2}(y_g-y_r,x_g-x_r)\). La velocidad angular se calcula con el error entre el angulo deseado y la orientacion actual del robot. Para evitar movimientos bruscos se saturan \(v\) y \(\omega\).

\section{Pseudocodigo}

\begin{lstlisting}
target = sim.getObject('/mannequin')
robot  = sim.getObject('/PioneerP3DX')

-- En cada paso de simulacion:
robotPos  = sim.getObjectPosition(robot, -1)
targetPos = sim.getObjectPosition(target, -1)
dx = targetPos[1] - robotPos[1]
dy = targetPos[2] - robotPos[2]
desiredHeading = math.atan2(dy, dx)
headingError = normalizeAngle(desiredHeading - robotYaw)
v = clamp(kv * math.sqrt(dx*dx + dy*dy), 0, vmax)
w = clamp(kw * headingError, -wmax, wmax)
setWheelSpeeds(v, w)
\end{lstlisting}

\section{Criterios de exito}

El robot debe avanzar hacia el objetivo, reducir velocidad cerca de la meta, evitar oscilaciones excesivas y detenerse dentro de una tolerancia definida. La entrega final debera incluir capturas de la escena y, si es posible, una grafica de distancia al objetivo.
